\vspace{-2pt}
\subsection{Properties of large autoregressive language models}
\vspace{-2pt}

\firstpara{Scaling laws} are found and studied in \textit{autoregressive} language models~\cite{scalinglaw,scalingar}, which describe a power-law relationship between the scale of model (or dataset, computation, \textit{etc.}) and the cross-entropy loss value on the test set.
Scaling laws allow us to directly predict the performance of a larger model from smaller ones \cite{gpt4}, thus guiding better resource allocation.
More pleasingly, they show that the performance of LLMs can scale well with the growth of model, data, and computation and never saturate, which is considered a key factor in the success of ~\cite{gpt3, llama1, llama2, opt, bloom, chinchilla}.
% Scaling laws also show within limited computational resources, larger models can be more data-efficient than smaller ones because they can reach a lower test loss with the same computational budget~\cite{scalinglaw,scalingar}.
% Some contributes the success of scaling laws to the autoregressive modeling' or something like this. Just give the evidence that scaling laws are mainly verified on AR rather Diffusion Transformers.}
The success brought by scaling laws has inspired the vision community to explore more similar methods for multimodality understanding and generation~\cite{llava,alayrac2022flamingo,visionllm,dong2023dreamllm,cm3leon_chameleon,emu_baai,chen2023internvl,dai2023emu_meta,jin2023unified,ge2023seed_llama,ge2024seedx,tian2024mminterleaved,wang2024git}.

% Propelled by the principles of scaling laws in neural language models~\cite{scalinglaw}, several Large Language Models~\cite{gpt3, opt, bloom, chinchilla, llama1, llama2} have been proposed, embodying the principle that increasing the scale of models tends to yield enhanced performance outcomes. GPT~\cite{gpt2, gpt3}, predicated on a transformer decoder architecture, undergoes generative pre-training and scales the model size to an unprecedented 175B parameters. 
% LLama~\cite{llama1, llama2} release a collection of pre-trained and fine-tuned large language models (LLMs) ranging in scale from 7 billion to 70 billion parameters. 
% The manifest efficacy of scaling laws applied to language models has proffered a glimpse into the promising potential of upscale in visual models~\cite{lvm}. 






\para{Zero-shot generalization.} Zero-shot generalization~\cite{multitask_zeroshot} refers to the ability of a model, particularly a Large Language Model, to perform tasks that it has not been explicitly trained on.
% In light of scaling laws~\cite{scalinglaw} and the emergent capabilities~\cite{emergent} of larger models, present-day Large Language Models can now attain sterling performance on a multitude of downstream tasks via zero-shot learning, evincing excellent generalization abilities. These capabilities render such models suitable as foundation models~\cite{foundationmodels}, obviating the need for excessive task-specific architectural tailoring for downstream applications~\cite{bert}. 
Within the realm of the computer vision, there is a burgeoning interest in the zero-shot and in-context learning abilities of foundation models, CLIP~\cite{clip}, SAM~\cite{sam}, Dinov2~\cite{dinov2}. Innovations like Painter~\cite{painter} and LVM~\cite{lvm} extend visual prompters~\cite{visualprompttuning1,visualprompttuning2} to achieve in-context learning in vision.
% thereby facilitating the generalization to downstream unseen tasks.


\vspace{-3pt}
\subsection{Visual generation}
\vspace{-3pt}
\firstpara{Raster-scan autoregressive models} for visual generation necessitate the encoding of 2D images into 1D token sequences.
% Language models rely on tokenizer like BPE~\cite{bpe} to tokenize text data.
Early endeavors~\cite{igpt,van2016pixelcnn} have shown the ability to generate RGB (or grouped) pixels in the standard row-by-row, raster-scan manner.
\cite{reed2017mspixelcnn} extends \cite{van2016pixelcnn} by using multiple independent trainable networks to do super-resolution repeatedly.
VQGAN~\cite{vqgan} advances \cite{igpt,van2016pixelcnn} by doing autoregressive learning in the latent space of VQVAE~\cite{vqvae}.
It employs GPT-2 decoder-only transformer to generate tokens in the raster-scan order, like how ViT~\cite{vit} serializes 2D images into 1D patches.
VQVAE-2~\cite{vqvae2} and RQ-Transformer~\cite{rq} also follow this raster-scan manner but use extra scales or stacked codes.
Parti~\cite{parti}, based on the architecture of ViT-VQGAN~\cite{vit-vqgan}, scales the transformer to 20B parameters and works well in text-to-image synthesis.

\para{Masked-prediction model.} MaskGIT~\cite{maskgit} employs a VQ autoencoder and a masked prediction transformer similar to BERT~\cite{bert, beit, mae} to generate VQ tokens through a greedy algorithm.
MagViT~\cite{magvit} adapts this approach to videos, and MagViT-2~\cite{magvit2} enhances~\cite{maskgit,magvit} by introducing an improved VQVAE for both images and videos.
MUSE~\cite{muse} further scales MaskGIT to 3B parameters.
% setting new benchmarks in text-to-image synthesis.

\para{Diffusion models}' progress has centered around improved learning or sampling~\cite{scorebased,ddim, dpm-solver,dpmpp,bao2022analytic}, guidance~\cite{cfg, glide}, latent learning~\cite{ldm}, and architectures~\cite{cdm, dit, imagen, raphael}.
% The Denoising Diffusion Probabilistic Model~\cite{ddpm} leverages diffusion processes within the pixel space to predict noise, while the Cascade Diffusion Model~\cite{cascade-diffusion} employs a cascading mechanism of pixel diffusion models to iteratively refine the generated images. 
DiT and U-ViT~\cite{dit,bao2023all} replaces or integrates the U-Net with transformer, and inspires recent image~\cite{chen2023pixart,chen2024pixart_sigma} or video synthesis systems~\cite{bar2024lumiere,gupta2023photorealistic} including Stable Diffusion 3.0~\cite{stable-diffusion3}, SORA~\cite{sora}, and Vidu~\cite{bao2024vidu}.
